\subsection{Increasing and Decreasing Area}

From previous chapters, we have studied how to find out the increasing or decreasing area of a function. 

\subsubsection{Recall}
\noindent For an Interval $I$:
$$\forall x, y \in I, \forall z \in \R, (x < z < y, \implies z \in I)$$
Given a function $f: I \to R$
$$\forall z \in I, f'(z) > 0 \implies f(x) \text{ is increasing in interval } I \text{ on }\R$$
$$\forall z \in I, f'(z) < 0 \implies f(x) \text{ is decreasing in interval } I \text{ on }\R$$

\subsubsection{Summary}
\begin{itemize}
    \item Intervals where the function increases/decreases are found by first finding where the derivative value is 0 or undefined
    \item Test for the signs on the derivative function to the left/right of these x values
    \item Regions of increased by positive slope values: regions of decrease are indicated by negative slope values. 
\end{itemize}

\begin{example}
    For the given function below, find the intervals of increase/decrease. 
    $$h(x) = \frac{
        x - 2
    }{
        x^2 + 3x - 4
    }$$
    Solution:\\
    To start of, let's get the derivative of h(x) first
    \begin{align*}
        h'(x) &= \frac{
            (x^2 + 3x - 4) - (x - 2)(2x + 3)
        }{(x^2 + 3x - 4)^2} \\
        &= \frac{x^2 + 3x - 4 - 2x^2 + x + 6}{
            (x^2 + 3x - 4)^2
        }\\
        &= \frac{
            -x^2 + 4x + 2
        }{
            (x^2 + 3x - 4)^2
        } \\
        &= -\frac{x^2 - 4x - 2}{(x + 4)^2 (x - 1)^2}
    \end{align*}
    Find the critical number of the function. \\

    When $h'(x) = 0, x = 2 \pm \sqrt{6}$\\

    When $h'(x) = \text{undefined}, x = -4 \text{ or } x = 1$\\

    \noindent Then, we can do the line test. Follow the advanced function strategy. 
\end{example}